\title{Introduction to regression}
\author{Paul Schrimpf}
\institute{UBC \\ Economics 326}
\date{\today}

\begin{document}

\frame{\titlepage}
%\setcounter{tocdepth}{2}

\begin{frame}[shrink]
  \frametitle{Review of last week}
  \begin{itemize}
  \item Expectations and conditional expectations    
    \begin{itemize}
    \item Linear
    \item Iterated expectations
    \end{itemize}
  \item Asymptotics --- using large sample distribution to approximate
    finite sample distribution of estimators
    \begin{itemize}
    \item LLN: sample moments converge in probability to population moments, 
      \[ \underbrace{\frac{1}{n} \sum_{i=1}^n g(x_i)}_{\text{sample
          moment}} \inprob \underbrace{\Er[g(x)]}_{\text{population
          moment}} \]
    \item CLT: centered and scaled sample moments converge in
      distribution to population moments
      \[ \underbrace{\sqrt{n}}_{\text{``scaling''}} \left( \frac{1}{n}
        \sum_{i=1}^n g(x_i) \underbrace{-
          \Er[g(x)}_{\text{``centering''}}] \right) \indist N\left(0,
        \var(g(x))\right) \]
    \item Using CLT to calculate p-values
    \end{itemize}
  \end{itemize}
\end{frame}

%%%%%%%%%%%%%%%%%%%%%%%%%%%%%%%%%%%%%%%%%%%%%%%%%%%%%%%%%%%%%%%%%%%%%%
\part{Definition and interpretation of regression}
\frame{\partpage}

\begin{frame}
  \tableofcontents  
\end{frame}

\begin{frame}
  \frametitle{References}
  \begin{itemize}
  \item Main texts:
    \begin{itemize}
    \item \cite{ap2014} chapter 2
    \item \cite{w2013} chapter 2    
    \item \cite{sw2009} chapter 4-5
    \end{itemize}
    % \item These slides draw heavily from \cite{angrist2009}, which is
    %   unfortunately not published
  \item More advanced:
    \begin{itemize}
    \item \cite{ap2009} chapter 3 up to and including section 3.1.2
      (pages 27-40)
    \item \cite{bierens2010sreg}            
    \item \cite{abbring2001} chapter 3
    \item \cite{baltagi2002} chapter 3
    \item \cite{linton2017} chapters 16-20, 22
    \end{itemize}
  \item More introductory:
    \begin{itemize}
    \item \cite{dbc2012} chapter 7
    \end{itemize}
  \end{itemize}
\end{frame}

The most important things to understand are contained in
\cite{ap2014}. The most useful references for details are \cite{w2013} or
\cite{sw2009}.  \cite{ap2009} is also very nice, but a bit more
dense. \cite{dbc2012} is a simple introduction to regression with many
examples and not much math. PDFs of \cite{linton2017} and
\cite{baltagi2002} are available from UBC library. They are more
technical and difficult than Wooldridge, but I think would still be
useful.  \cite{bierens2010sreg} has many of the proofs that we will go
through, but the typesetting is not great.  \cite{abbring2001} moves
quickly and uses matrix notation. If you have not taken linear algebra
or you are not comfortable with matrix multiplication and inversion
(which you do not need to know and will not be covered in this
course), then \cite{abbring2001} may not be useful for you.

\section{Motivation}

\begin{frame}
  \frametitle{General problem}
  \begin{itemize}
  \item Often interested in relationship between two (or more)
    variables, e.g.\
    \begin{itemize}
    \item Wages and education
    \item Minimum wage and unemployment
    \item Price, quantity, and product characterics
    \end{itemize}
  \item Usually have:
    \begin{enumerate}
    \item Variable to be explained (\alert{dependent variable})
    \item \alert{Explanatory variable(s)} or \alert{independent
        variables} or \alert{covariates} 
    \end{enumerate}
    \begin{tabular}{c c} {\textbf{Dependent}} & {\textbf{Independent}}
      \\
      Wage & Education \\
      Unemployment & Minimum wage \\
      Quantity & Price  and product characteristics \\
      $Y$ & $X$ 
    \end{tabular}
  \item For now agnostic about causality, but $\Er[Y|X]$ usually is
    not causal
  \end{itemize}
\end{frame}

\begin{frame}
  \frametitle{Example: Growth and GDP}

  \includegraphics[width=\figwidth,height=\figheight]{growGDP-scatter}
     
\end{frame}

\begin{frame}\frametitle{Years of schooling in 1960 and growth}

  \includegraphics[width=\figwidth,height=\figheight]{growSchool-scatter}   

\end{frame}


\section{Conditional expectation function}

\begin{frame}[allowframebreaks]
  \frametitle{Conditional expectation function}
  \begin{itemize}
  \item One way to describe relation between two variables is a
    function, 
    \[ Y = h(X) \]
  \item Most relationships in data are not deterministic, so look at
    average relationship,
    \begin{align*}
      Y = & \underbrace{\Er[Y|X]}_{\equiv h(X)} + \underbrace{\left(Y
          - \Er[Y|X]\right)}_{\equiv \epsilon}  \\  
      = & \Er[Y|X] + \epsilon
    \end{align*}
  \item Note that $\Er[\epsilon] = 0$ (by definition of $\epsilon$ and
    iterated expectations)
  \item $\Er[Y|X]$ can be any function, in particular, it need not be
    linear
  \item Unrestricted $\Er[Y|X]$ hard to work with
    \begin{itemize} 
    \item Hard to estimate
    \item Hard to communicate if $X$ a vector (cannot draw graphs) 
    \end{itemize}
  \item Instead use linear regression
    \begin{itemize}
    \item Easier to estimate and communicate
    \item Tight connection to $\Er[Y|X]$
    \end{itemize}
  \end{itemize}
\end{frame}  

\section{Population regression} 
\begin{frame}[allowframebreaks]\frametitle{Population regression}
  \begin{itemize}
  \item The bivariate \alert{population regression} of $Y$ on $X$ is 
    \[ (\beta_0, \beta_1) = \argmin_{b_0,b_1} \Er[(Y - b_0 - b_1
    X)^2 ] \] 
    i.e. $\beta_0$ and $\beta_1$ are the slope and intercept that
    minimize the expected square error of $Y - (\beta_0 + \beta_1 X)$ 
  \item Calculating $\beta_0$ and $\beta_1$:
    \begin{itemize}
    \item First order conditions:
      \begin{align}
        [b_0]: 0 = & \frac{\partial}{\partial b_0} \Er[(Y - b_0 - b_1
        X)^2 ] \notag \\
         = & \Er\left[\frac{\partial}{\partial b_0}  (Y - b_0 - b_1 X)^2 \right]
         \notag \\
        = & \Er\left[ -2(Y - \beta_0 - \beta_1 X) \right] \label{b0}
      \end{align}
      and
      \begin{align}
        [b_1]:  0 = & \frac{\partial}{\partial b_1} \Er[(Y - b_0 - b_1
        X)^2 ] \notag \\
         = & \Er\left[\frac{\partial}{\partial b_1}  (Y - b_0 - b_1 X)^2 \right]
        \notag \\
        = & \Er\left[ -2(Y - \beta_0 - \beta_1 X)X \right] \label{b1}
      \end{align}
    \item (\ref{b0}) rearranged gives $\beta_0 = \Er[Y] - \beta_1
      \Er[X]$
    \item Substituting into (\ref{b1})
      \begin{align*}
        0 = & \Er\left[ X(-Y + \Er[Y] - \beta_1 \Er[X] + \beta_1 X) \right] \\
        = & \Er\left[X(-Y + \Er[Y])\right] + \beta_1\Er\left[X(X -
          \Er[X])\right] \\
        = & -\cov(X,Y) + \beta_1 \var(X) \\
        \beta_1 = & \frac{\cov(X,Y)}{\var(X)}
      \end{align*}
    \end{itemize}
  \item \alert{$\beta_1 = \frac{\cov(X,Y)}{\var(X)}$, $\beta_0 = \Er[Y] - \beta_1
      \Er[X]$}
  \end{itemize}
\end{frame}

\begin{frame}\frametitle{Population regression approximates
    $\Er[Y|X]$}
  \begin{lemma} \label{lem:regapprox}
    The population regression is the minimal mean square error linear
    approximation to the conditional expectation function, i.e.\
    {\footnotesize{
        \begin{align*}
      \underbrace{\argmin_{b_0, b_1} \Er\left[\left(Y-(b_0 + b_1
          X) \right)^2\right]}_{\text{population regression}} =
      \argmin_{b_0, b_1} \underbrace{\Er_X\left[\left(\Er[Y|X] - (b_0 + b_1
          X)\right)^2\right]}_{\text{MSE of linear approximation to $\Er[Y|X]$}}
    \end{align*}
  }}
  \end{lemma}
  
  \begin{corollary} 
    If $\Er[Y|X] = c + m X$, then the population regression of $Y$ on
    $X$ equals $\Er[Y|X]$, i.e. $\beta_0 = c$ and $\beta_1 = m$
  \end{corollary}
\end{frame}

\begin{frame}[shrink]\frametitle{Proof}
  \begin{proof}
    \begin{itemize}
    \item Let $b^*_0, b^*_1$ be minimizers of MSE of approximation to $\Er[Y|X]$
    \item Same steps as in population regression formula gives
      \[ 0 = \Er\left[ -2(\Er[Y|X] - b^*_0 - b^*_1 X) \right]  \]
      and
      \[ 0 = \Er\left[ -2(\Er[Y|X] - b^*_0 - b^*_1 X)X \right]  \]
    \item Rearranging and combining,
      \begin{align*}
        b_0^* = \Er[\Er[Y|X]] - b_1^* \Er[X] = \Er[Y] - b_1^* \Er[X]
      \end{align*}
      and
      \begin{align*}
        0 = & \Er\left[ X(-\Er[Y|X] + \Er[Y] + b^*_1 \Er[X] - b^*_1 X) \right] \\
        = & \Er\left[X(-\Er[Y|X] + \Er[Y])\right] + b^*_1\Er\left[X(X -
          \Er[X])\right] \\
        = & -\cov(X,Y) + b^*_1 \var(X) \\
        b^*_1 = & \frac{\cov(X,Y)}{\var(X)}      
      \end{align*}
    \end{itemize}
  \end{proof}
\end{frame}

\subsection{Interpretation}
\begin{frame} \frametitle{Regression interpretation}
  \begin{itemize}
  \item Regression $=$ best linear approximation to $\Er[Y|X]$
  \item $\beta_0 \approx \Er[Y|X=0]$
  \item $\beta_1 \approx \frac{d}{dx} \Er[Y|X] \approx$ change in average
    $Y$ per unit change in $X$
  \item Not necessarily a causal relationship (usually not)
  \item Always can be viewed as description of data
  \end{itemize}
\end{frame}

\begin{frame}\frametitle{Regression with binary $X$}
  \begin{columns}[c] % align columns
    \begin{column}{.48\textwidth}
      \begin{itemize}
      \item Suppose $X$ is binary (i.e. can only be $0$ or $1$)
      \item We know $\beta_0 + \beta_1 X = $ best linear approximation to
        $\Er[Y|X]$ 
      \item $X$ only takes two values, so can draw line connecting
        $\Er[Y|X=0]$ and $\Er[Y|X=1]$, so $\beta_0 + \beta_1 X = \Er[Y|X]$
        \begin{itemize}
        \item $\beta_0 = \Er[Y|X=0]$
        \item $\beta_0 + \beta_1 = \Er[Y|X=1]$
        \end{itemize}
      \end{itemize}
    \end{column}
    \begin{column}{.48\textwidth}
      \includegraphics[width=\textwidth]{binary-regression}
    \end{column}
  \end{columns}
\end{frame}

% \subsection{Properties}

\section{Sample regression}
\begin{frame} \frametitle{Sample regression}
  \begin{itemize}
  \item Have sample of observations: $\{ (y_i, x_i) \}_{i=1}^n$ 
  \item The \alert{sample regression} (or when unambiguous just
    ``regression'') of $Y$ on $X$ is 
    \[ (\hat{\beta}_0, \hat{\beta}_1) = \argmin_{b_0,b_1}
    \frac{1}{n}\sum_{i=1}^n (y_i - b_0 - b_1 x_i)^2 \] 
    i.e. $\hat{\beta}_0$ and $\hat{\beta}_1$ are the slope and intercept that
    minimize the sum of squared errors, $(y_i - (\hat{\beta}_0 + \hat{\beta}_1 x_i))^2$ 
    \begin{itemize}
    \item Same as population regression but with sample average
      instead of expectation
    \end{itemize}
  \item Same calculation as for population regression would show
    \[ \hat{\beta}_1 = \frac{\widehat{\cov}(X,Y)}{\widehat{\var(X)}} =
    \frac{\frac{1}{n} \sum_{i=1}^n (x_i -\bar{x})(y_i - \bar{y})}
    {\frac{1}{n} \sum_{i=1}^n (x_i - \bar{x})^2} \]
    and 
    \[ \hat{\beta}_0 = \bar{y} - \hat{\beta}_1 \bar{x} \]
  \end{itemize}
\end{frame}

Since $\hat{\beta}_1$ and $\hat{\beta}_0$ come from minimizing a sum
of squares, they are called the ordinary least squares estimates, or
OLS for short. 

The formulas for $\hat{\beta}_1$ and $\hat{\beta}_0$ come from the
first order conditions.  These estimators minimize the sum of squared
differences between the regression line and the observed $y_i$, 
\[ (\hat{\beta}_0, \hat{\beta}_1) = \argmin_{b_0,b_1}
\frac{1}{n}\sum_{i=1}^n (y_i - b_0 - b_1 x_i)^2. \] 
The first order condition for $\hat{\beta}_0$ is:
\begin{align*}
  0 = & \frac{1}{n} \sumin 2(y_i - \hat{\beta}_0 - \hat{\beta}_1 x_i ) 
\end{align*}
which can be rearranged to get
\begin{align*}
  \hat{\beta}_0 = & \frac{1}{n} \left(\sumin y_i - \hat{\beta}_1 x_i
  \right) \\
  \hat{\beta}_0 = & \bar{y} - \hat{\beta}_1 \bar{x}.
\end{align*}
The first order condition for $\hat{\beta}_1$ is: 
\begin{align*}
  0 = & \frac{1}{n} \sumin 2x_i (y_i - \hat{\beta}_0 - \hat{\beta}_1 x_i ).
\end{align*}
Substituting in the previous expression for $\hat{\beta}_0$ gives
\begin{align*}
  0 = & \frac{1}{n} \sumin 2 x_i (y_i - \bar{y} + \hat{\beta}_1 \bar{x} - \hat{\beta}_1 x_i ). 
\end{align*}
Rearranging to solve for $\hat{\beta}_1$:
\begin{align*}
  \hat{\beta}_1 \sumin x_i(x_i - \bar{x}) = & \sumin x_i (y_i -
  \bar{y}) \\
  \hat{\beta}_1 = & \frac{\sumin x_i (y_i - \bar{y})}{ x_i (x_i -
    \bar{x}) } = \frac{\sumin (x_i - \bar{x}) (y_i - \bar{y})}{ (x_i -
    \bar{x})^2 }.
\end{align*}

\begin{frame}
  \frametitle{Sample regression}
  \begin{itemize}
  \item Sample regression is an estimator for the population
    regression
  \item Given an estimator we should ask:
    \begin{itemize}
    \item Unbiased?
    \item Variance?
    \item Consistent?
    \item Asymptotically normal?
    \end{itemize}
  \item We will address these questions in the next week or two  
  \end{itemize}
\end{frame}

% \begin{frame}[shrink]
%   \frametitle{True linear model approach to regression}
%   \begin{itemize}
%   \item Most authors (like Wooldridge) introduce regression by
%     starting with a linear model for $Y$:
%     \[ Y = \alpha_0 + \alpha_1 X + \underbrace{U}_{\text{unobserved}} \] 
%     with $\Er[U X] = 0$
%     \begin{itemize}
%     \item Perspective: model $=$ true description of data generating
%       process 
%     \end{itemize}
%   \item Implications:
%     \begin{itemize}
%     \item Population regression coefficients $=$ model coefficients,
%       i.e. $\beta_0 = \alpha_0$ and $\beta_1 = \alpha_1$
%     \item Conditional expectation function is linear 
%       \[ \Er[Y|X] = \alpha_0 + \alpha_1 X \]
%     \item $\beta_1=\alpha_1$ has causal interpretation as long as we
%       believe $\Er[U X] = 0$
%     \item Easier to discuss causality
%     \item Easier to derive statistical properties
%     \end{itemize}
%   \end{itemize}
% \end{frame}

% \begin{frame}[shrink]
%   \frametitle{Problems with true linear models}
%   \begin{itemize}
%   \item Usually do not believe models are linear, e.g.\ no economic
%     theory that says the following should be linear:
%     \begin{enumerate}
%     \item $\log (wage_i) = \beta_0 + \beta_1 (educ_i) + \epsilon_i$
%     \item $\log q_t = \beta_0 + \beta_1 \log p_t + \epsilon_t$
%     \end{enumerate}
%   \item Usually do not believe error terms are uncorrelated with
%     covariates, e.g.\
%     \begin{enumerate}
%     \item Need $\Er[u_i educ_i]=0$, but $u_i$ probably includes IQ,
%       propensity to work hard, etc, which should be correlated with
%       education 
%     \item Is it supposed to be demand or supply? Either case, changes
%       in $q_t$ from $\epsilon_t$ generally also change equilibrium $p_t$, so
%       $\Er[\log p_t \epsilon_t ] \neq 0$
%     \end{enumerate}
%   \item Viewing regression as best linear approximation to $\Er[Y|X]$
%     makes it clear what regression tells you about the data even if
%     the true model is not linear and does not have $\Er[XU]=0$
%   \end{itemize}
% \end{frame}


\section{Regression in R}

%% slide about true linear model vs approx CEF

\begin{frame}[fragile,shrink] \frametitle{Regression in R}
  \begin{lstlisting}
require(datasets) ## some datasets included with R
stateDF <- data.frame(state.x77)
summary(stateDF) ## summary statistics of data

## Sample regression function
regress <- function(y, x) {
  beta <- vector(length=2)
  beta[2] <- cov(x,y)/var(x)
  beta[1] <- mean(y) - beta[2]*mean(x)
  return(beta)
}

## Regress life expectancy on income
beta <- regress(stateDF[,"Life.Exp"],stateDF$Income)
beta

## builtin regression 
lm(Life.Exp ~ Income, data=stateDF)
## more detailed output
summary(lm(Life.Exp ~ Income, data=stateDF))
  \end{lstlisting}
  \url{https://bitbucket.org/paulschrimpf/econ326/src/master/notes/03/regress.R?at=master}
\end{frame}
%$
%%%%%%%%%%%%%%%%%%%%%%%%%%%%%%%%%%%%%%%%%%%%%%%%%%%%%%%%%%%%%%%%%%%%%%
\clearpage
\part{Properties of regression}

\frame{\partpage}

\section{Fitted value and residuals}

These algebraic identities about fitted values and residuals are
things that we will use repeatedly later. I would not recommend
spending time trying to memorize these. The important ones will come
up repeatedly and you will remember them without any special
effort. The first time we use these identities, we will go through how
to get them again. We may even go through them yet again the second
and third time we use them. Eventually we will use some of these
identities so often that you will either be able to quickly derive
them or just remember them.

\begin{frame}[allowframebreaks]
  \frametitle{Fitted values and residuals}
  \begin{itemize}
  \item \alert{Fitted values}: 
    \[ \hat{y}_i = \hat{\beta}_0 +
    \hat{\beta}_1 x_i \]
  \item \alert{Residuals}: 
    \[\hat{\epsilon}_i = y_i - \hat{\beta}_0 -
    \hat{\beta}_1 x_i = y_i  - \hat{y}_i \] 
    \[ y_i = \hat{y}_i + \hat{\epsilon}_i \]
  \item Sample mean of residuals $=0$
    \begin{itemize}
    \item First order condition for $\hat{\beta}_0$,
      \begin{align*}
        0 = &\avg (y_i - \hat{\beta}_0 - \hat{\beta}_1
        x_i) \\ 
        0 = & \avg \hat{\epsilon}_i 
      \end{align*}
    \end{itemize}
  \item Sample covariance of $x$ and $\hat{\epsilon}$ $= 0$
    \begin{itemize}
    \item First order condition for $\hat{\beta}_1$,
      \begin{align*}
        0 = &\avg (y_i - \hat{\beta}_0 - \hat{\beta}_1
        x_i)x_i \\ 
        0 = & \avg \hat{\epsilon}_i x_i
      \end{align*}
    \end{itemize}
  \item Sample mean of $\hat{y}_i = \bar{y} = \hat{\beta}_0 +
    \hat{\beta}_1 \bar{x}$ 
    \begin{align*}
      \avg y_i = & \avg \hat{y}_i + \hat{\epsilon}_i \\
      = & \avg \hat{y}_i \\
      = & \avg  \hat{\beta}_0 + \hat{\beta}_1 x_i \\
      = & \hat{\beta}_0 + \hat{\beta}_1 \bar{x}
    \end{align*}
  \item Sample covariance of $y$ and $\hat{\epsilon}$ = sample
    variance of $\hat{\epsilon}$:
    \begin{align*}
      \avg y_i (\hat{\epsilon}_i - \bar{\hat{\epsilon}}) = & \avg y_i
      \hat{\epsilon}_i \\
      = & \avg (\hat{\beta}_0 + \hat{\beta}_1 x_i +
      \hat{\epsilon}_i)\hat{\epsilon}_i \\
      = & \hat{\beta}_0 \avg \hat{\epsilon}_i + \beta_1 \avg x_i
      \hat{\epsilon}_i + \avg \hat{\epsilon}_i^2 \\
      = & \avg \hat{\epsilon}_i^2 
    \end{align*}
  \end{itemize}
\end{frame}

\begin{frame}
  \frametitle{$R^2$}
  \begin{itemize}
  \item Decompose $y_i$
    \[ y_i = \hat{y}_i + \hat{\epsilon}_i \]
  \item Total sum of squares $=$ explained sum of squares $+$ sum of
    squared residuals
    \[ \underbrace{\avg (y_i - \bar{y})^2}_{SST} =
    \underbrace{\avg (\hat{y}_i - \bar{y})^2}_{SSE} +
    \underbrace{\avg \hat{\epsilon}_i^2}_{SSR} \]
  \item \alert{R-squared}: fraction of sample variation in $y$ that is
    explained by $x$
    \[ R^2 = \frac{SSE}{SST} = 1 - \frac{SSR}{SST} =o
    \widehat{\corr}(y,\hat{y}) \]
    \begin{itemize}
    \item $0 \leq R^2 \leq 1$
    \item If all data on regression line, then $R^2 = 1$
    \item Magnitude of $R^2$ does not have direct bearing on economic
      importance of a regression
    \end{itemize}
  \end{itemize}
\end{frame}

\section{Statistical properties}

\subsection{Unbiased}
\begin{frame}
  \frametitle{Unbiased}
  \begin{itemize}
  \item $\Er[\hat{\beta}] = ?$
  \item Assume:
    \begin{enumerate}
    \item[SLR.1]\label{s1} (linear model) $ y_i = \beta_0 + \beta_1
      x_i + \epsilon_i $
    \item[SLR.2]\label{s2} (independence) $\{(x_i,y_i)\}_{i=1}^n$ is independent random
      sample
    \item[SLR.3]\label{s3} (rank condition) $\widehat{\var}(X) > 0$
    \item[SLR.4]\label{s4} (exogeneity) $\Er[\epsilon | X] = 0$
    \end{enumerate}
  \item Then, $\Er[\hat{\beta}_1] = \beta_1$ and $\Er[\hat{\beta}_0] = \beta_0$
  \end{itemize}
\end{frame}

It is more important to understand the meaning of these four
assumptions (discussed below) than the proof that regression is
unbiased.  
\begin{proof}[Regression is unbiased]
  We need to calculate $\Er[\hat{\beta}]$. First, substitute in the
  formula for $\hat{\beta}$.
  \begin{align*}
    \Er[\hat{\beta}_1] = & \Er\left[ \frac{\sumin (x_i - \bar{x}) y_i
      }{\sumin (x_i - \bar{x})x } \right] 
  \end{align*}
  Next, substitute in the model for $y_i$, $y_i = \beta_0 + \beta_1
  x_i + \epsilon_i$, 
  \begin{align*}
    \Er[\hat{\beta}_1] = & \Er\left[ \frac{\sumin (x_i - \bar{x}) (\beta_0 + \beta_1
  x_i + \epsilon_i)
      }{\sumin (x_i - \bar{x})x } \right] \\
    & \text{ rearrange } \\
    = & \Er\left[ \frac{ \overbrace{\sumin x_i - \bar{x}}^{=0}}{\sumin
        (x_i - \bar{x})x } \beta_0 + \overbrace{\left(\frac{\sumin (x_i - \bar{x})
            x_i}{\sumin (x_i - \bar{x})x }\right)}^{=1} \beta_1 +
      \frac{\sumin (x_i - \bar{x})\epsilon_i} {\sumin (x_i - \bar{x})x }
    \right] \\
    & \text{ use linearity of expectation } \\ 
    = & \beta_1 + \Er\left[ \frac{\sumin (x_i - \bar{x})\epsilon_i} {\sumin (x_i - \bar{x})x }
    \right] \\
    & \text{ use iterated expectations } \\
    = & \beta_1 + \Er_X\left[ \Er_{\epsilon|X}\left[\left. \frac{\sumin (x_i -
            \bar{x})\epsilon_i} {\sumin (x_i - \bar{x})x} \right\vert x_1,
        x_2, ..., x_n \right]\right] \\
    & \text{ conditional on $x_1, ..., x_n$, $x_i$ is constant} \\
    = & \beta_1 + \Er_X\left[ \frac{\sumin (x_i -
        \bar{x})\Er[\epsilon_i|x_1, ..., x_n]} {\sumin (x_i -
        \bar{x})x} \right] \\
    & \text{ independent observations implies $\Er[\epsilon_i|x_1, ...,
      x_n] = \Er[\epsilon_i | x_i]$ } \\
    = & \beta_1 + \Er_X\left[ \frac{\sumin (x_i -
        \bar{x})\Er[\epsilon_i|x_i]} {\sumin (x_i - \bar{x})x} \right] \\
    & \text{ exogeneity assumption says that $\Er[\epsilon_i | x_i] =
      0$.} \\ 
    = & \beta_1
  \end{align*}
\end{proof}
Note that the first few steps of the above proof showed that 
\begin{align*}
  \hat{\beta}_1 = \beta_1 + \frac{\sumin (x_i - \bar{x})
    \epsilon_i}{\sumin (x_i - \bar{x})^2}. 
\end{align*}
This is a very useful expression that can be used as a starting point
for calculating the variance of $\hat{\beta}_1$, and thinking about
what happens if exogeneity fails and $\Er[\epsilon|x] \neq 0$.

\subsection{Variance}

\begin{frame}
  \frametitle{Variance}
  \begin{itemize}
  \item $\var(\hat{\beta})?$
  \item Assume SLR.1-4 and
    \begin{enumerate}
    \item[SLR.5]\label{s5} (homoskedasticity) $\var(\epsilon | X) =
      \sigma^2$
    \end{enumerate}
  \item Then, 
    \[ \var(\hat{\beta}_1 | \{x_i\}_{i=1}^n ) =
    \frac{\sigma^2}{\sum_{i=1}^n (x_i - \bar{x})^2} \]
    and 
    \[ \var(\hat{\beta}_0 | \{x_i\}_{i=1}^n ) =
    \frac{\sigma^2 \avg x_i^2 }{\sum_{i=1}^n (x_i - \bar{x})^2} \]
  \end{itemize}
\end{frame}

As in the proof that regression is unbiased, 
\begin{align*}
  \hat{\beta}_1 = \beta_1 + \frac{\sumin (x_i - \bar{x})
    \epsilon_i}{\sumin (x_i - \bar{x})^2}. 
\end{align*}
We now want to take the variance of this expression. Before doing so,
it will be useful to review some properties of the variance of a sum
of random variables. 
\begin{lemma}
  Let $a$, $b$, and $c$ be constants, and $Z$ and $W$ be random
  variables. Then,
  \[ \var(a + bZ + cW) = b^2 \var(Z) + c^2 \var(W) + 2 bc
  \cov(Z,W). \]
\end{lemma}
\begin{proof}
  We can prove this using the definition of variance. 
  \begin{align*}
    \var(a + bZ + cW) = & \Er\left[ \left(a + bZ + cW - \Er[a + bZ + cW]\right)^2 \right] \\
    = & \Er\left[ \left(a + bZ + cW - a - b\Er[Z] - c\Er[W]\right)^2
    \right] \\
    = & \Er\left[\left(b(Z - \Er[Z]) + c(W - \Er[W]) \right)^2 \right]
    \\
    = & \Er\left[ b^2 (Z - \Er[Z])^2 + 2bc (Z - \Er[Z])(W - \Er[W]) +
      c^2 (W - \Er[W])^2 \right] \\
    = & b^2\Er\left[  (Z - \Er[Z])^2\right] + 2bc \Er\left[ (Z - \Er[Z])(W - \Er[W]) \right]+
    c^2 \Er\left[ (W - \Er[W])^2 \right] \\
    = & b^2 \var(Z) + c^2 \var(W) + 2 bc \cov(Z,W)
  \end{align*}
\end{proof}
Generalizing the above to the sum of more than two random variables,
we have
\begin{corollary}
  Let $a_1, ..., a_n$ be constants, and $Z_1, ..., Z_n$ be random
  variables, then,
  \[
  \var\left( \sumin a_i Z_i \right)  = \sumin \sum_{j=1}^n a_i a_j
  \cov(Z_i, Z_j) 
  \]
  Futhermore, if $Z_i$ and $Z_j$ are independent (or just
  uncorrelated) for $i \neq j$, then
  \[
  \var\left( \sumin a_i Z_i \right)  = \sumin a_i^2 \var(Z_i)
  \]
\end{corollary}
We can apply this corollary to 
\begin{align*}
  \var(\hat{\beta}_1|x) = & \var\left(\left. \beta_1 + \frac{\sumin (x_i - \bar{x})
        \epsilon_i}{\sumin (x_i - \bar{x})^2} \right\vert x \right)  \\
  = & \var\left(\left. \sumin \underbrace{\frac{x_i - \bar{x}}{\sumin (x_i -
          \bar{x})^2}}_{a_i} \underbrace{\epsilon_i}_{Z_i} \right\vert x
  \right)
  & \text{ using the corollary } \\
  = & \sumin \sum_{j=1}^n \frac{x_i - \bar{x}}{\sumin (x_i -
    \bar{x})^2} \frac{x_j - \bar{x}}{\sumin (x_i -
    \bar{x})^2} \cov(\epsilon_i, \epsilon_j | x)
  & \text{ independence } \\
  = & \sumin \left(\frac{x_i - \bar{x}}{\sumin x_i - \bar{x}}
  \right)^2 \var(\epsilon_i | x) 
  & \text{ homoskedasticity} \\
  = & \sumin \frac{(x_i - \bar{x})^2}{\left( \sumin (x_i - \bar{x})^2
    \right)^2} \sigma_\epsilon^2 \\
  = & \frac{\sigma_\epsilon^2}{\sumin (x_i - \bar{x})^2}.
\end{align*}

\subsection{Distribution}
\begin{frame}
  \frametitle{Distribution with normal errors}
  \begin{itemize}
  \item Assume SLR.1-SLR.5 and
    \begin{enumerate}
    \item[SLR.6] (normality) $\epsilon_i|x_i \sim N(0,\sigma^2$)
    \end{enumerate}
  \item Then $Y|X \sim N(\beta_0 + \beta_1 X, \sigma^2)$, and
    \[ \hat{\beta}_1|\{x_i\}_{i=1}^n \sim N\left(\beta_1, \frac{\sigma^2}{\sum_{i=1}^n
        (x_i - \bar{x})^2} \right) \]
  \item Even without assuming normality, the central limit theorem
    implies $\hat{\beta}$ is asymptotically normal (details in a later
    lecture)
  \end{itemize}
\end{frame}

An important property of normal random variables is that if $Z$ and
$W$ are independent, and $Z \sim N(\mu_z, \sigma_z^2)$ and $W \sim
N(\mu_w, \sigma_w^2)$, then 
\[ a + bZ + cW \sim N(a + b \mu_z + c \mu_w, b^2 \sigma^2_z + c^2
\sigma^2_w). \] If we assume that $\epsilon_i$ is normally distributed
conditional on $x$, then since $\hat{\beta}_1$ is just a sum of the
$\epsilon_i$,\footnote{Specifically, $ \hat{\beta}_1 = \beta_1 +
  \frac{\sumin (x_i - \bar{x}) \epsilon_i}{\sumin (x_i -
    \bar{x})^2}$.}
$\hat{\beta}_1$ will also be normally distributed. 

\begin{frame} \frametitle{Summary}
  \begin{itemize}
  \item Simple linear regression model assumptions:
    \begin{enumerate}
    \item[SLR.1] (linear model) $ Y_i = \beta_0 + \beta_1 x_i + \epsilon_i $
    \item[SLR.2] (independence) $\{(x_i,y_i)\}_{i=1}^n$ is independent random
      sample
    \item[SLR.3] (rank condition) $\widehat{\var}(X) > 0$
    \item[SLR.4] (exogeneity) $\Er[\epsilon | X] = 0$
    \item[SLR.5] (homoskedasticity) $\var(\epsilon | X) =
      \sigma^2$
    \item[SLR.6] (normality) $\epsilon_i|x_i \sim N(0,\sigma^2)$
    \end{enumerate}
  \item $\hat{\beta}$ unbiased if SLR.1-SLR.4
  \item If also SLR.5, then $\var(\hat{\beta}_1 | \{x_i\}_{i=1}^n ) =
    \frac{\sigma^2}{\sum_{i=1}^n (x_i - \bar{x})^2}$ 
  \item If also SLR.6, then $\hat{\beta}_1|\{x_i\}_{i=1}^n \sim
    N\left(\beta_1, \frac{\sigma^2}{\sum_{i=1}^n (x_i - \bar{x})^2}
    \right)$
  \end{itemize}
\end{frame}

\subsection{Discussion of assumptions}
\begin{frame}[allowframebreaks]
  \frametitle{Discussion of assumptions}
  \begin{itemize}
  \item[SLR.1] Having a linear model makes it easier to state the other
    assumptions, but we could instead start by saying let $\beta_1 =
    \frac{\cov(X,Y)}{\var(X)}$ and $\beta_0 = \Er[Y] - \beta_1 \Er[X]$
    be the population regression coefficients and define $\epsilon_i =
    y_i - \beta_0 - \beta_1 x_i$
  \end{itemize}
\end{frame}

To say whether an estimator is unbiased, we first have to define what
parameter we want to estimate. Assuming that there is a linear model
defines the parameter we want to estimate. This linear model could be
population regression, in which case, $\beta_1 = \frac{\cov(X,Y)
}{\var(X)}$, and by construction we must have
$\Er[X\epsilon]=0$. 

However, the linear model may also be motivated by economic
theory. For example, consider a Cobb-Douglass production function with
only one input, labor,
\[ Y = A L^\alpha, \] 
where $Y$ is output, $L$ is labor, and $A$ is productivity. If we take logs, then
\[ \log Y = \log A + \alpha \log L. \]
If we rearrange slightly, we get something that looks just like a
linear regression model, 
\[ \underbrace{\log Y_i}_{y_i} = \underbrace{\Er[\log A]}_{\beta_0}
+ \underbrace{\alpha}_{\beta_1} \underbrace{\log L_i}_{x_i} +
\underbrace{(\log{A}_i - \Er[\log A])}_{\epsilon_i}. \]
If this is the model we want to estimate, then $\epsilon_i$ is not the
error term in the population regression. Instead $\epsilon_i$ is the
difference between the log productivity of firm $i$ and average log
productivity. It is unlikely that this $\epsilon_i$ would be
uncorrelated with $\log L_i$. More productive firms generally choose
to use more inputs, so we should suspect that $\epsilon_i$ and $\log
L_i$ are positively correlated.

\begin{frame}[allowframebreaks]
  \frametitle{Discussion of assumptions}
  \begin{itemize}
  \item[SLR.2] Independent observations is a good assumption for data
    from a simple random sample
    \begin{itemize}
    \item Common situations where it fails in
      economics are when we have a time series of observations,
      e.g.\ $\{(x_t,y_t)\}_{t=1}^n$ could be unemployment and GDP of
      Canada for many different years; and clustering, e.g.\ the data
      could be students test scores and hours studying and our sample
      consists of randomly chosen courses or schools---students in the
      same course would not be independent, but across different courses
      they might be.
    \item Still have $\Er[\hat{\beta}_1] = \beta_1$ with
      non-independent observations as long as $\Er[\epsilon_i | x_1, ..., x_n] = 0$
    \item The variance of $\hat{\beta}_1$ will change with
      non-independent observations
    \item
      \href{https://bitbucket.org/paulschrimpf/econ326/src/master/notes/03/dependent.R?at=master}
      {Simulation code}
    \end{itemize}
  \end{itemize}
\end{frame}

Independence says that knowing the values of $x_1$ and $y_1$ tells you
nothing about the distribution of $x_2$ and $y_2$ (or any other
observation). When we have cross-sectional data, this assumption
usually makes sense. In economics, we sometimes deal with time-series
data, where $(x_1,y_1)$ would be the observation of something at time
$1$ and $(x_2,y_2)$ is the observation that same thing at time $2$. In
this case, independence is unlikely to hold. Another common situation
is panel data, where we observe a sample of individuals over time, so
$(x_{it}, y_{it})$ would be what we observe from individual $i$ at
time $t$. Again, it is unlikely that these observations would be
independent over time.  Later in the course, we will talk about how to
deal with non-independent observations.

\begin{frame}[allowframebreaks]
  \frametitle{Discussion of assumptions}
  \begin{itemize}
  \item[SLR.3] If $\widehat{\var}(X) = 0$, then 
    $\hat{\beta}_1$ involves dividing by $0$
    \begin{itemize}
    \item If there is no variation in $X$, then we cannot see how
      $Y$ is related to $X$
    \end{itemize}
  \end{itemize}
\end{frame}

\begin{frame}[allowframebreaks]
  \frametitle{Discussion of assumptions}
  \begin{itemize}
  \item[SLR.4] To think about mean independence of $\epsilon$ from $x$
    we should have a model motivating the regression
  \end{itemize}
  \begin{columns}[c] % align columns
    \begin{column}{.48\textwidth}
      \begin{itemize}
      \item If the model we want is just a population regression,
        then automatically $\Er[\epsilon X] = 0$, and $\Er[\epsilon|X] =
        0$ if the conditional expectation function is linear; if
        conditional expectation nonlinear maybe still a useful
        approximation       
      \end{itemize}
    \end{column}
    \begin{column}{.48\textwidth}
      \includegraphics[width=\textwidth]{nonlinear} \\
      \href{https://bitbucket.org/paulschrimpf/econ326/src/master/notes/03/nonlinear.R?at=master}
      {Code}
    \end{column}
  \end{columns}
\end{frame}

\begin{frame}[allowframebreaks]
  \frametitle{Discussion of assumptions}
  \begin{itemize}
  \item[SLR.4] To think about mean independence of $\epsilon$ from $x$
    we should have a model motivating the regression
  \end{itemize}
  \begin{columns}[c] % align columns
    \begin{column}{.48\textwidth}
      \begin{itemize}
      \item If the model we want is anything else, then maybe
        $\Er[\epsilon X] \neq 0$ (and $\Er[\epsilon | X] \neq 0$), e.g.\ 
        \begin{itemize}
        \item Demand curve 
          \[ p_i = \beta_0 + \beta_1 q_i + \epsilon_i \]
          $\epsilon_i = $ everything that affects price other than
          quantity. $q_i$ determined in equilibrium implies
          $\Er[\epsilon_i | q_i] \neq 0$ 
        \item $\Er[\hat{\beta}_1] \neq \beta_1$ and $\hat{\beta}_1$ does
          not tell us what we want
        \end{itemize}
      \end{itemize}
    \end{column}
    \begin{column}{.48\textwidth}
      \includegraphics[width=\textwidth]{endogenous} \\
      \href{https://bitbucket.org/paulschrimpf/econ326/src/master/notes/03/endogenous.R?at=master}
      {Code}
    \end{column}
  \end{columns}
\end{frame}
  

Exogeneity is the most important assumption underlying regression. In
fact, estimating any economic model using any method will involve some
kind of exogeneity assumption. By this, we mean that every estimation
method requires assuming some error term is either completely
independent of some observable ($F_{\epsilon|x}(\epsilon|x) =
F_\epsilon(\epsilon)$), mean independent of some observable
($\Er[\epsilon|X] = 0$), or at least uncorrelated with an observable
($\Er[\epsilon X] = 0$). Much of what separates good empirical work in
economics from bad is how plausible are the exogeneity
assumptions. Often, economic theory can help us decide whether or
not an exogeneity assumption is plausible. Consider the 
production function example from earlier, 
\[ \underbrace{\log Y_i}_{y_i} = \underbrace{\Er[\log A]}_{\beta_0}
+ \underbrace{\alpha}_{\beta_1} \underbrace{\log L_i}_{x_i} +
\underbrace{(\log{A}_i - \Er[\log A])}_{\epsilon_i}. \]
To think about whether mean independence of the error term,
$\Er[\underbrace{(\log{A} - \Er[\log A])}|\log L]=0$, makes sense in
this model, we should think about how $L$ is determined. The firm
chooses how much labor to use. Suppose the firm faces output price $p$
and wage $w$. If the firm chooses $L$ knowing its productivity, then
the firm solves,
\begin{align*}
   \max_L p A L^\alpha - wL 
\end{align*}
The first order condition is 
\begin{align*}
  p A \alpha L^{\alpha -1} - w = & 0.
\end{align*}
If we solve for $A$, we get
\begin{align*}
  A = & \frac{w}{p \alpha } L^{1-\alpha} \\
  \log A = \log\left(\frac{w}{p \alpha }\right) + (1-\alpha) \log L
\end{align*}
so
\begin{align*}
  \Er[ \log A | \log L ] =  (1-\alpha) \log L +
  \Er\left[\left . \log\left(\frac{w}{p \alpha }\right) \right| \log L
  \right] 
\end{align*}
This will not be a constant unless $\alpha = 1$ and $\frac{w}{p}$ is
mean independent of $\log L$. Both of these are unlikely to
hold. Unless $\Er[ \log A | \log L ] $ is constant,
$\Er[\underbrace{(\log{A} - \Er[\log A])}|\log L] \neq 0$.  Therefore,
exogeneity is not a good assumption in this model, and regression will
not give an unbiased estimate of the production function. 

\begin{frame}[allowframebreaks]
  \frametitle{Discussion of assumptions}
  \begin{itemize}
  \item[SLR.5] Homoskedasticity: variance of $\epsilon$ does not
    depend on $X$ \\
    \begin{tabular}{cc} 
      \alert{Homoskedastic} & \alert{Heteroskedastic} \\ 
      \includegraphics[width=0.4\textwidth]{homoskedastic} &
      \includegraphics[width=0.4\textwidth]{heteroskedastic} 
    \end{tabular} \\
    \href{https://bitbucket.org/paulschrimpf/econ326/src/master/notes/03/skedastic.R?at=master}
    {Code}
    \begin{itemize}
    \item \alert{Heteroskedasticity} is when $\var(\epsilon|X)$ varies
      with $X$
    \item If there is heteroskedasticity, the variance of
      $\hat{\beta}_1$ is different, but we can fix it 
    \item ``robust standard errors'' / ``heteroscedasticity-consistent
      (HC) standard errors'' / ``Eicker–Huber–White standard errors''
    \end{itemize}
  \end{itemize}
\end{frame}

Homoskedasticity is a strong assumption that is usually not very
plausible. Therefore, in practice, economists almost always calculate
heteroscedasticity-robust standard errors.

\begin{frame}[allowframebreaks]
  \frametitle{Discussion of assumptions}
  \begin{itemize}
  \item[SLR.6] If $\epsilon_i|x_i \sim N$, then $\hat{\beta}_1 \sim N$
    \begin{itemize}
    \item What if $\epsilon_i$ not normally distributed?
    \item We will see that $\hat{\beta}_1$ still asymptotically normal
    \item
      \href{https://bitbucket.org/paulschrimpf/econ326/src/master/notes/03/reg-clt.R?at=master}
      {Simulation}
    \end{itemize}
  \end{itemize}
\end{frame}


\begin{frame}
  \frametitle{Discussion of assumptions}

  \includegraphics[width=\figwidth,height=\figheight,keepaspectratio]{beta1} 

\end{frame}

\begin{frame}
  \frametitle{Discussion of assumptions}

  \includegraphics[width=\figwidth,height=\figheight,keepaspectratio]{zbeta1} 

\end{frame}

\section{Examples}
\begin{frame}\frametitle{Example: smoking and cancer}
  \begin{itemize}
  \item Data on per capita number of cigarettes sold and death rates
    per thousand from cancer for U.S. states in 1960
  \item \url{http://lib.stat.cmu.edu/DASL/Datafiles/cigcancerdat.html}
  \item Death rates from: lung cancer, kidney cancer, bladder cancer,
    and leukemia
    \href{https://bitbucket.org/paulschrimpf/econ326/src/master/notes/03/smoking.R?at=master} {Code}
  \end{itemize}
\end{frame}

\begin{frame}\frametitle{Smoking and lung cancer}

  \includegraphics[width=\figwidth,height=\figheight]{cigLung}   

\end{frame}

\begin{frame}\frametitle{Smoking and kidney cancer}

  \includegraphics[width=\figwidth,height=\figheight]{cigKidney}   

\end{frame}

\begin{frame}\frametitle{Smoking and bladder cancer}

  \includegraphics[width=\figwidth,height=\figheight]{cigBladder}   

\end{frame}

\begin{frame}\frametitle{Smoking and leukemia}

  \includegraphics[width=\figwidth,height=\figheight]{cigLeukemia}   

\end{frame}

\begin{frame}\frametitle{Example: convergence in growth}
  \begin{itemize}
  \item Data on average growth rate from 1960-1995 for 65 countries
    along with GDP in 1960, average years of schooling in 1960, and
    other variables
  \item From
    \url{http://wps.aw.com/aw_stock_ie_2/50/13016/3332253.cw/index.html},
    originally used in \cite{beck2000}
  \item Question: has there been in convergence, i.e.\ did poorer
    countries in 1960 grow faster and catch-up?    
  \item \href{https://bitbucket.org/paulschrimpf/econ326/src/master/notes/03/growth.R?at=master} {Code}
  \end{itemize}
\end{frame}

\begin{frame}\frametitle{GDP in 1960 and growth}

  \includegraphics[width=\figwidth,height=\figheight]{growGDP}   

\end{frame}

\begin{frame}\frametitle{Years of schooling in 1960 and growth}

  \includegraphics[width=\figwidth,height=\figheight]{growSchool}   

\end{frame}

\begin{frame}
  \begin{itemize}
  \item Things look different 1995-2014
  \item
    \href{https://bitbucket.org/paulschrimpf/econ326/src/master/notes/03/growth-updated.R?at=master}
    {Code to download and recreate results using updated growth data
      through 2014 from the World Bank}
  \end{itemize}
\end{frame}

%%%%%%%%%%%%%%%%%%%%%%%%%%%%%%%%%%%%%%%%%%%%%%%%%%%%%%%%%%%%%%%%%%%%%%
\section{Inference}
\begin{frame}[allowframebreaks]
  \frametitle{Inference with normal errors}
  \begin{itemize}
  \item Regression estimates depend on samples, which are random, so
    the regression estimates are random
    \begin{itemize}
    \item Some regressions will randomly look ``interesting'' due to
      chance 
    \end{itemize}
  \item Logic of hypothesis testing: figure out probability of getting
    an interesting regression estimate due solely to change
  \item Null hypothesis, $H_0:$ the regression is uninteresting,
    usually $\beta_1 = 0$
  \item With assumptions SR.1-SR.6 and under $H_0: \beta_1 =
    \beta_1^\ast$, we know
    \[ \hat{\beta} \sim N\left(\beta_1^\ast,
      \frac{\sigma^2_\epsilon}{\sum_{i=1}^n (x_i - \bar{x})^2}
    \right) \]
    or equivalently,
    \[ t \equiv \frac{\hat{\beta} - \beta_1^\ast}{\sigma_\epsilon /
      \sqrt{\sum_{i=1}^n (x_i - \bar{x})^2} } \sim N(0,1) \]
  \item P-value: the probability of getting a regression estimate as
    or more ``interesting'' than the one we have
    \begin{itemize}
    \item As or more interesting $=$ as far or further away from
      $\beta_1^\ast$
    \item If we are only interested when $\hat{\beta}_1$ is on one
      side of $\beta_1^\ast$, then we have a one sided alternative,
      e.g.\ $H_a: \beta_1 > \beta_1^\ast$
    \item If we are equally interested in either direction, then $H_a:
      \beta_1 \neq \beta_1^\ast$
    \end{itemize}
  \end{itemize}
\end{frame}

\begin{frame} 

  \includegraphics[width=\figwidth,height=\figheight]{oneSided} 

\end{frame}

\begin{frame} 

  \includegraphics[width=\figwidth,height=\figheight]{twoSided} 

\end{frame}
  
\begin{frame}[allowframebreaks]
  \frametitle{Inference with normal errors}
  \begin{itemize}
  \item One-sided p-value: $p = \Phi(-\abs{t}) = 1-\Phi(\abs{t})$ \\
  \item Two-sided p-value: $p = 2\Phi(-\abs{t}) = 2(1-\Phi(\abs{t}))$
  \item Interpretation:
    \begin{itemize}
    \item The probability of getting an estimate as strange as the one
      we have if the null hypothesis is true.
    \item It is \emph{not} about the probability of $\beta_1$ being
      any particular value. $\beta_1$ is not a random variable. It is
      some unknown number. The data is what is random. In particular,
      the p-value is \emph{not} the probability that that $H_0$ is false
      given the data.
    \end{itemize}
  \item Hypothesis testing: we must make a decision (usually reject or
    fail to reject $H_0$)
    \begin{itemize}
    \item Choose significance level $\alpha$ (usually 0.05 or 0.10)
    \item Construct procedure such that if $H_0$ is true, we will
      incorrectly reject with probability $\alpha$
    \item Reject null if p-value less than $\alpha$
    \end{itemize}
  \end{itemize}
\end{frame}

\subsection{Examples (continued)}
\begin{frame}\frametitle{Smoking and cancer}

  \input{fig/smoking}

\end{frame}

\begin{frame}\frametitle{Growth and GDP}

  \input{fig/growth}

\end{frame}

\begin{frame}\frametitle{Caution: multiple testing}
  \begin{itemize}
  \item We just looked at 6 regressions, if $H_0: \beta_1 = 0$ is true
    in all of them the probability that correctly fail to reject all 6
    null hypotheses with a 5\% test is $0.95^6 = 0.74$ (assuming the 6
    tests are independent)
  \item A quarter of the time if we look at 6 regressions, we will
    randomly find at least significant relationship; if we look at 14
    regressions the probability that we incorrectly reject a null is
    more than 0.5
  \end{itemize}
\end{frame}

\begin{frame}\frametitle{Caution: economic significance $\neq$
    statistical significance} 
\end{frame}
%%%%%%%%%%%%%%%%%%%%%%%%%%%%%%%%%%%%%%%%%%%%%%%%%%%%%%%%%%%%%%%%%%%%%% 
\subsection{Estimating $\sigma_\epsilon^2$}

\begin{frame}[allowframebreaks] 
  \frametitle{Estimating $\sigma_\epsilon^2$}
  \begin{itemize}
  \item Recall that $\var(\hat{\beta}|x_1,...,x_n) =
    \frac{\sigma^2_\epsilon}{\sumin (x_i - \bar{x})^2} =
    \frac{\sigma_\epsilon^2}{n \widehat{\var}(x)}$
  \item $\sigma_\epsilon^2$ unknown
  \item We estimate $\sigma_\epsilon^2$ using the residuals,
    \[ \hat{\sigma}_\epsilon^2 = \frac{1}{n-2} \sumin
    \underbrace{\hat{\epsilon}_i^2}_{=(y_i - \hat{\beta}_0 -
      \hat{\beta}_1 x_i)^2} \]
  \item If SLR.1-SLR.5, $\Er[\hat{\sigma}_\epsilon^2 ] = \sigma_\epsilon^2$ 
  \item Using $\frac{1}{n-2}$ instead of $\frac{1}{n}$ makes
    $\hat{\sigma}_\epsilon^2$ unbiased
    \begin{itemize}
    \item $\hat{\epsilon}_i$ depends on 2 estimated parameters,
      $\hat{\beta}_0$ and $\hat{\beta}_1$, so only $n-2$ degrees of
      freedom 
    \end{itemize}    
  \item Estimate $\var(\hat{\beta}_1|x_1,...,x_n)$ by 
    \[ \widehat{\var}(\hat{\beta}_1|x_1,...,x_n) =
    \frac{\hat{\sigma}^2_\epsilon}{\sumin (x_i - \bar{x})^2} =
    \frac{\frac{1}{n-2} \sumin \hat{\epsilon}_i^2}{\sumin (x_i -
      \bar{x})^2} \]
  \item Standard error of $\hat{\beta}_1$ is
    $\sqrt{\widehat{\var}(\hat{\beta}_1|x_1,...,x_n)}$ 
  \item If SLR.1-SLR.6, t-statistic with estimated
    $\widehat{\var}(\hat{\beta}_1|x_1,...,x_n)$ has a $t(n-2)$
    distribution instead of $N(0,1)$
    \[ t = \frac{\hat{\beta}_1 -
      \beta_1}{\sqrt{\widehat{\var}(\hat{\beta}_1|x_1, ..., x_n)}} \sim
    t(n-2) \]    
  \end{itemize}
\end{frame}
%%%%%%%%%%%%%%%%%%%%%%%%%%%%%%%%%%%%%%%%%%%%%%%%%%%%%%%%%%%%%%%%%%%%%%
\subsection{Confidence intervals}

\begin{frame}[allowframebreaks]\frametitle{Confidence intervals}
  \begin{itemize}
  \item $\hat{\beta}_1$ is random
  \item $\widehat{\var}(\hat{\beta}_1)$, p-values, and hypthesis tests
    are ways of expressing how random is $\hat{\beta}_1$ 
  \item Confidence intervals are another
  \item A $1-\alpha$ confidence interval, $CI_{1-\alpha} =
    [LB_{1-\alpha}, UB_{1-\alpha}]$ is an interval estimator for
    $\beta_1$ such that 
    \[ \Pr(\beta_1 \in CI_{1-\alpha} = 1-\alpha) \]
    ($CI_{1-\alpha}$ is random; $\beta_1$ is not)
  \item Recall: if SLR.1-SLR.6, then 
    \[ \hat{\beta}_1 \sim N\left(\beta_1,\var(\hat{\beta}_1)\right) \]
    \begin{itemize}
    \item Implies
      \begin{align*} 
        \Pr\left(\hat{\beta}_1 < \beta_1 + \sqrt{\var(\hat{\beta}_1)}
          \Phi^{-1}(\alpha/2)\right) = & \alpha/2 \\
        \Pr\left(\hat{\beta}_1 - \sqrt{\var(\hat{\beta}_1)}
          \Phi^{-1}(\alpha/2) < \beta_1)\right) = & \alpha/2 \\
      \end{align*}
      and 
      \begin{align*} 
        \Pr\left(\hat{\beta}_1 > \beta_1 + \sqrt{\var(\hat{\beta}_1)}
          \Phi^{-1}(1-\alpha/2) \right) = & \alpha/2 \\
        \Pr\left(\hat{\beta}_1 - \sqrt{\var(\hat{\beta}_1)}
          \Phi^{-1}(1-\alpha/2) > \beta_1 \right) = & \alpha/2
      \end{align*}
      so 
      \begin{align*}
        \Pr & \begin{pmatrix} \hat{\beta}_1 + \sqrt{\var(\hat{\beta}_1)}
          \Phi^{-1}(\alpha/2) < \beta_1 \\ \beta_1 < \hat{\beta} +
          \sqrt{\var(\hat{\beta}_1)} \Phi^{-1}(1-\alpha/2)
        \end{pmatrix}
        = \\
        = & 1- \Pr\left(\hat{\beta}_1 + \sqrt{\var(\hat{\beta}_1)}
          \Phi^{-1}(\alpha/2) < \beta_1\right) - \\
         & - 
        \Pr\left(\hat{\beta}_1 + \sqrt{\var(\hat{\beta}_1)}
          \Phi^{-1}(1-\alpha/2) > \beta_1\right) \\
        = & 1 - \alpha
      \end{align*}
    \item For $\alpha = 0.05$, $\Phi^{-1}(0.025) \approx -1.96$,
      $\Phi^{-1}(0.975) \approx 1.96$
    \item For $\alpha = 0.1$, $\Phi^{-1}(0.05) \approx -1.64$
    \end{itemize}
  \end{itemize}
\end{frame}

\begin{frame}
  \frametitle{Confidence intervals}

  \includegraphics[width=\figwidth,height=\figheight]{ci}

\end{frame}

\begin{frame}\frametitle{Confidence intervals}
  \begin{itemize}
  \item $1-\alpha$ confidence interval
    \[ \hat{\beta}_1 \pm \sqrt{\var(\hat{\beta}_1)}
    \Phi^{-1}(\alpha/2) \]
  \item With estimated $\hat{\sigma}_\epsilon^2$, use t distribution
    instead of normal
    \[ \hat{\beta}_1 \pm \sqrt{\widehat{\var}(\hat{\beta}_1)}
    F_{t,n-2}^{-1}(\alpha/2) \]
    $F_{t,n-2}^{-1} = $ inverse CDF of $t(n-2)$ distribution \\
    \begin{tabular}{c|cccccc}
      \multicolumn{7}{c}{{\textbf{$F_{t,n-2}(1-\alpha/2)$}}} \\
      & \multicolumn{6}{c}{$n-2$} \\
      $\alpha/2$ & 5 & 10 & 20 & 50 & 100 & $\infty$ \\ \hline
      0.025 & 2.57 & 2.23 & 2.09 & 2.01 & 1.98 & 1.96 \\
      0.05 &  2.02 & 1.81 & 1.72 & 1.68 & 1.66 & 1.64 \\ \hline
    \end{tabular}
  \end{itemize}  
\end{frame}

\begin{frame}\frametitle{Example: GDP in 1960 and growth}

  \includegraphics[width=\figwidth,height=\figheight]{growGDP-ci}   

\end{frame}

\begin{frame}\frametitle{Example: Years of schooling in 1960 and
    growth}

  \includegraphics[width=\figwidth,height=\figheight]{growSchool-ci}   

\end{frame}

%%%%%%%%%%%%%%%%%%%%%%%%%%%%%%%%%%%%%%%%%%%%%%%%%%%%%%%%%%%%%%%%%%%%%% 
\section{Efficiency}
\begin{frame}[shrink]
  \frametitle{Gauss-Markov theorem}
  \begin{itemize}
  \item The sample regression estimator,
    \begin{align*} (\hat{\beta}_0, \hat{\beta}_1) = & \argmin \sumin (y_i - b_0 -
      b_1 x_i)^2 \\
      \hat{\beta}_1 = & \frac{\sumin (x_i - \bar{x}) y_i}{\sumin(x_i -
        \bar{x})^2}
    \end{align*}
    also called \alert{O}rdinary \alert{L}east \alert{S}quares (OLS)
    is not the only unbiased estimator of 
    \[ y_i = \beta_0 + \beta_1 x_i + \epsilon_i \]
  \item \alert{Gauss-Markov theorem:} if SLR.1-SLR.5, then OLS is the
    \alert{B}est \alert{L}inear \alert{U}nbiased \alert{E}stimator
    \begin{itemize}
    \item Linear means linear in $y$, $\hat{\beta}_1 = \sumin c_i y_i$
      with $c_i = \frac{(x_i - \bar{x})}{\sumin (x_i - \bar{x})^2}$
    \item Unbiased means $\Er[\hat{\beta}_1] = \beta_1$
    \item Best means that among all linear unbiased estimators, OLS
      has the smallest variance
    \end{itemize}
  \end{itemize}
\end{frame}

\begin{frame}\frametitle{Proof: setup}
  \begin{itemize}
  \item Let $\tilde{\beta}_1$ be a linear unbiased estimator of
    $\beta_1$
    \begin{itemize}
    \item Linear: $\tilde{\beta}_1 = \sumin c_i y_i$
    \item Unbiased: $\Er[\tilde{\beta}_1|x_1, ...,x_n] = \beta_1$ (for all possible
      $\beta_0,\beta_1$)    
    \end{itemize}
  \item We will show that 
    \[ \var(\tilde{\beta}_1 | x_1, ... , x_n) \geq \var(\hat{\beta}_1
    | x_1, ..., x_n) \]
  \end{itemize}  
\end{frame}

The $c_i$'s in the definition of $\tilde{\beta}_1$ are allowed to be
nonlinear functions of $x_i$. For example OLS uses $c_i = \frac{(x_i -
  \bar{x})}{\sum_{j=1}^n (x_j - \bar{x})^2}$. 

\begin{frame}\frametitle{Proof: outline}
  \begin{enumerate}
  \item Show that $\sumin c_i = 0$ and $\sumin c_i x_i = 1$
  \item Show $\cov(\tilde{\beta}_1, \hat{\beta}_1|x_1,...,x_n) =
    \var(\hat{\beta}_1|x_1, ..., x_n)$
  \item Show $\var(\tilde{\beta}_1 | x_1, ..., x_n) \geq
    \var(\hat{\beta}_1 | x_1, ..., x_n)$
  \item Show $\var(\tilde{\beta}_1 | x_1, ..., x_n) =
    \var(\hat{\beta}_1 | x_1, ..., x_n)$ only if $\tilde{\beta}_1 =
    \hat{\beta}_1$ 
  \end{enumerate}
  \footnotetext{We will go over the proof in class. See
    \href{http://faculty.arts.ubc.ca/vmarmer/econ326/326_05_gauss_markov_handout.pdf}
    {Marmer's slides} or \cite{w2013} 3A for details}
\end{frame}

\begin{theorem}[Gauss-Markov] 
  If SLR.1-SLR.5, then OLS is the \alert{B}est \alert{L}inear
  \alert{U}nbiased \alert{E}stimator. In other words if 
  \[ \tilde{\beta}_1 = \sum_{i=1}^n c_i y_i \]
  is another linear (in $y$) unbiased ($\Er[\tilde{\beta}_1|x_1, ..., x_n] = \beta_1$)
  estimator, then 
  \[ \var(\tilde{\beta}_1 | x_1, ..., x_n) \geq \var(\hat{\beta}_1 |
  x_1, ..., x_n) \]
  where the inequality is strict unless $\tilde{\beta}_1 = \hat{\beta}_1$.
\end{theorem}
\begin{proof}  
  \textbf{Step 1: show that $\sumin c_i = 0$ and $\sumin c_i x_i =
    1$.} We do this by using the fact that $\tilde{\beta}_1$ is
  unbiased. Note that 
  \begin{align*}
    \Er[\tilde{\beta}_1|x_1, ..., x_n] = & \Er[\sumin c_i y_i |x_1,
                                           ..., x_n] \\
    & \text{substitute in the model for $y$} \\
    = & \Er[\sumin c_i (\beta_0 + \beta_1 x_i + \epsilon_i)|x_1, ..., x_n ] \\
    & \text{use linearity of expectation} \\
    = & \beta_0 \sumin \Er[c_i | x_1, ..., x_n] + \beta_1 \sumin x_i
        \Er[c_i | x_1, ..., x_n]  + \sumin \Er[\epsilon_i c_i | x_1,
        ... , x_n] \\
    & \text{$c_i$ is constant given $x_1, ... , x_n$} \\
    = & \beta_0 \sumin c_i + \beta_1 \sumin x_i c_i \\    
  \end{align*}
  Hence, unbiasedness requires 
  \[ \beta_1 = \beta_0 \sumin c_i + \beta_1 \sumin x_i c_i.  \]
  An estimator being unbiased means that it is unbiased for any
  possible value of the true parameters. For the previous equation to
  hold for all possible $\beta_0$ and $\beta_1$, we need $\sumin c_i =
  0$ and $\sumin c_i x_i = 1$.
  
  \textbf{Step 2:  show $\cov(\tilde{\beta}_1, \hat{\beta}_1|x_1,...,x_n) =
    \var(\hat{\beta}_1|x_1, ..., x_n)$.} 
  Substituting in $y_i =
  \beta_0 + \beta_1 x_i + \epsilon_i$ let's us write $\tilde{\beta}_1$ as
  \begin{align*}
    \tilde{\beta}_1 = & \sumin c_i y_i \\
    = & \sumin c_i (\beta_0 + \beta_1 x_i  + \epsilon_i) \\
    = & \beta_0 (\sumin c_i) + \beta_1 (\sumin c_i x_i) + (\sumin c_i
        \epsilon_i) \\
    & \text{Using result of step 1} \\
    = & \beta_1 + \sumin c_i \epsilon_i
  \end{align*}
  Similarly for $\hat{\beta}_1$, 
  \begin{align*}
    \hat{\beta}_1 = & \beta_1 + \sumin \frac{(x_i
                      -\bar{x})}{\sum_{j=1}^n (x_j - \bar{x})^2}
                      \epsilon_i  
  \end{align*}
  Let $w_i = \frac{(x_i-\bar{x})}{\sum_{j=1}^n (x_j - \bar{x})^2}$ so
  that we can write 
  \[ \hat{\beta}_1 = \beta_1 + \sumin w_i \epsilon_i. \]
  Now let's calculate the covariance of $\tilde{\beta}_1$ and
  $\hat{\beta}_1$. 
  \begin{align*}
    \cov(\tilde{\beta}_1, \hat{\beta}_1 | x_1, ..., x_n) = 
    &
      \Er[(\tilde{\beta}_1 - \beta_1) (\hat{\beta}_1 - \beta_1) ] \\
    = & \Er\left[ \left. \left( \sumin c_i \epsilon_i \right) \left(\sumin
        w_i \epsilon_i \right)\right\vert x_1, ..., x_n \right] \\
  = & \sumin c_i w_i \Er[\epsilon_i^2 | x_1, ..., x_n ] + \sumin
        \sum_{j\neq i} c_i w_j \Er[\epsilon_i \epsilon_j  | x_1, ...,
        x_n] \\
    & \text{second term $= 0$ because of independent observations} \\
    = &  \sumin c_i w_i \Er[\epsilon_i^2 | x_1, ..., x_n] \\
    & \text{using homoskedasticity} \\
    = & \sigma^2 \sumin c_i w_i \\
    & \text{substituting } w_i = \frac{(x_i-\bar{x})}
      {\sum_{j=1}^n (x_j- \bar{x})^2} \\
    = & \sigma^2 \sumin c_i \frac{(x_i-\bar{x})}{\sum_{j=1}^n (x_j -
        \bar{x})^2} \\
    = & \frac{\sigma^2}{\sum_{j=1}^n (x_j -
        \bar{x})^2} \left((\sumin c_i x_i) - (\bar{x} \sumin c_i)
        \right) \\
    & \text{ using result of step 1} \\
    = & \frac{\sigma^2}{\sum_{j=1}^n (x_j -
        \bar{x})^2} \left(1 - \bar{x} 0\right)\\
    = & \var(\hat{\beta}_1 | x_1, ..., x_n)
  \end{align*}
  
  \textbf{Step 3:  show $\var(\tilde{\beta}_1 | x_1, ..., x_n) \geq
    \var(\hat{\beta}_1 | x_1, ..., x_n)$.} Consider
  $\var(\tilde{\beta}_1 - \hat{\beta}_1|x_1, ..., x_n)$. We know that this is $\geq
  0$ because it is a variance. Also, 
  \begin{align*}
    \var(\tilde{\beta}_1 - \hat{\beta}_1|x_1, ..., x_n) = &
    \var(\tilde{\beta}_1 | x_1, ..., x_n) +     \var(\hat{\beta}_1 |
    x_1, ..., x_n) - 2\cov(\tilde{\beta}_1,\hat{\beta}_1 | x_1, ...,
    x_n) \\
    & \text{Using result of step 2} \\
    = & \var(\tilde{\beta}_1 | x_1, ..., x_n) -\var(\hat{\beta}_1
        |x_1, ..., x_n) 
  \end{align*}
  Hence, 
  \begin{align*}
    \var(\tilde{\beta}_1 - \hat{\beta}_1|x_1, ..., x_n) =
    \var(\tilde{\beta}_1 | x_1, ..., x_n) -\var(\hat{\beta}_1   |x_1,
    ..., x_n)   & \geq 0 \\
    \var(\tilde{\beta}_1 | x_1, ..., x_n) & \geq \var(\hat{\beta}_1 |
                                            x_1, ..., x_n)
  \end{align*}
  
  \textbf{Step 4: show
    $\var(\tilde{\beta}_1 | x_1, ..., x_n) = \var(\hat{\beta}_1 | x_1,
    ..., x_n)$
    only if $\tilde{\beta}_1 = \hat{\beta}_1$} This step is only
  needed for the ``where the inequality is strict only if
  $\tilde{\beta}_1 = \hat{\beta}_1$ part of the theorem. If
  $\var(\tilde{\beta}_1 | x_1, ..., x_n) = \var(\hat{\beta}_1 | x_1,
  ..., x_n)$, then
  \[ \var(\tilde{\beta}_1 - \hat{\beta}_1|x_1, ..., x_n) =
  \var(\tilde{\beta}_1 | x_1, ..., x_n) -\var(\hat{\beta}_1   |x_1,
  ..., x_n)   = 0, \]
  so $\tilde{\beta}_1 - \hat{\beta}_1$ is constant. Since both are
  unbiased, the constant difference must be zero. Thus,
  $\tilde{\beta}_1 = \hat{\beta}_1$. 
\end{proof}



% \begin{frame}
%  \tableofcontents  
%\end{frame}

\begin{frame}[allowframebreaks]
  \frametitle{References}
\bibliographystyle{jpe}
\bibliography{../326}
\end{frame}
  

\end{document}